$latexcv-preamble.tex()$

%-------------------------------------------------------------------------------
%	TWO COLUMN
%
%	The body set in two balanced columns with compact stacked entry headlines to
%	suit the narrower measure. The name is a filled bar opening the left column,
%	the contact line runs across the head of the page, and a second filled bar
%	closes the foot.
%
%	Upstream builds its two columns by hand, as a run of paired 0.485\textwidth
%	minipages with each section placed into one of them. That only balances for
%	the content upstream happens to ship, and a Markdown body is a single
%	stream with nowhere to say which column a section belongs in. `multicol` is
%	used here instead, so the columns balance themselves and the content may
%	run to any length.
%-------------------------------------------------------------------------------

% The bottom margin is only as deep as the footer bar plus a little air above
% it: the bar bleeds off the foot of the paper, so what lies below the body is
% the bar itself and not a margin. Upstream gets the same effect with
% `bottom=-.6cm`, running the text block past the paper edge, because its bar is
% the last thing in the document rather than the page foot.
$if(geometry)$
\geometry{$for(geometry)$$geometry$$sep$,$endfor$}
$else$
\geometry{top=1.5cm, bottom=24pt, left=1.25cm, right=1.25cm}
$endif$

% The contact line is set in the running head, so the page must actually
% reserve room for one. The shared preamble leaves \headheight at upstream's
% -5pt, which is only safe while the head is empty. Applied after the block
% above, so it holds whether or not the document supplied its own `geometry`.
\geometry{headheight=14pt, headsep=14pt}

% The distance from the text edge to the paper edge, for the elements that bleed
% out of the left column: the title bar, the photograph under it, and the footer.
% Measured at \begin{document}, once the geometry above has settled.
\newlength{\cvbleed}
\newlength{\cvbleedwidth}
\AtBeginDocument{\setlength{\cvbleed}{\dimexpr 1in+\oddsidemargin\relax}}

% The same measurement downwards: from the baseline fancyhdr sets the foot on to
% the bottom edge of the paper. \cvfooterbar lowers the bar by this, less the
% bar's own depth, so the bar ends on the paper edge whatever the page geometry
% turns out to be -- including a document that supplied its own.
\newlength{\cvfootbleed}
\AtBeginDocument{\setlength{\cvfootbleed}{%
  \dimexpr\paperheight-1in-\topmargin-\headheight-\headsep-\textheight-\footskip\relax}}

\setlength{\columnsep}{$if(columnsep)$$columnsep$$else$18pt$endif$}

% Space between an entry and the next, named because it is used in three
% places here.
\newcommand{\spread}{4pt}

%-------------------------------------------------------------------------------
%	HEADINGS
%-------------------------------------------------------------------------------
% \linewidth rather than \textwidth, so the bar fills the column it is set in
% rather than overrunning into the next.
\newcommand{\cvsection}[1]{%
  \vspace{4pt}
  \colorbox{bgcol}{\mystrut\makebox[\linewidth][l]{%
    \cvchevrons{sectcol}\hspace{4pt}%
    \textcolor{white}{\textbf{#1}}\hspace{4pt}}}\\
  \vspace{2pt}}

\newcommand{\cvsubsection}[1]{%
  \vspace{2pt}
  {\larrow{sectcol}\hspace{2pt}\textcolor{bgcol}{\textbf{#1}}}\par
  \vspace{2pt}}

%-------------------------------------------------------------------------------
%	ENTRIES
%-------------------------------------------------------------------------------
% Stacked rather than in columns: at half the page width there is no room for
% a date column beside the text. The date and context share the accent-coloured
% first line, with the title in bold beneath.
\newcommand{\cventryhead}[4]{%
  \textcolor{sectcol}{%
    \ifthenelse{\isempty{#1}}{}{#1}%
    \ifthenelse{\isempty{#3}}{}{\ifthenelse{\isempty{#1}}{}{ - }#3}%
    \ifthenelse{\isempty{#4}}{}{, \textit{#4}}}\\
  \textbf{#2}}

\newcommand{\vitaeentry}[5]{%
  \begin{flushleft}
    \cventryhead{#1}{#2}{#3}{#4}\\[-7pt]
    \textcolor{softcol}{\hrule}
    \vspace{\spread}
    #5
  \end{flushleft}
  \vspace{-\spread}}

\newcommand{\vitaebrief}[4]{%
  \begin{flushleft}
    \cventryhead{#1}{#2}{#3}{#4}
  \end{flushleft}
  \vspace{-\spread}}

% The measurement and the row itself are in the shared preamble: every design
% sets this listing the same way, and only the floor below differs.
\newcommand{\vitaeskill}[2]{\cvskillrow{#1}{#2}}

\newenvironment{vitaewhy}{%
  \begin{itemize}[label=\larrow{bgcol}, topsep=0pt, partopsep=0pt, itemsep=2pt,
                  parsep=0pt, leftmargin=1.4em, labelsep=0.3em]%
}{%
  \end{itemize}%
}

%-------------------------------------------------------------------------------
%	RUNNING HEAD
%-------------------------------------------------------------------------------
% Upstream sets name, position and contacts as one small centred line above the
% columns. That is what frees the title block to sit inside the left column, so
% the two go together.
%
% \cvfirstcontactfalse rather than \cvcontactsreset: the name and position have
% already been emitted, so the first contact wants a separator in front of it
% like every other. Setting the flag this way round also means a document with
% no contact fields at all emits nothing here rather than a dangling separator.
% It matters that this is set on every page -- the head is re-executed per page,
% and \cvcontacts otherwise remembers that it has already run.
% `$$` is how a Pandoc template writes a literal dollar, so the separator reaches
% LaTeX as the math-mode \cdot upstream uses rather than opening a template
% variable named \cdot.
\renewcommand{\cvcontactsep}{\hspace{5pt}\textcolor{softcol}{$$\cdot$$}\hspace{5pt}}
\chead{\small
  \textsc{$name$$if(surname)$ $surname$$endif$}%
  $if(position)$\cvcontactsep $position$$endif$%
  \cvfirstcontactfalse\cvcontacts}

%-------------------------------------------------------------------------------
%	FOOTER
%-------------------------------------------------------------------------------
% A filled bar across the foot of the page, running the full width of the paper.
%
% Upstream fakes one at the very end of its document with \vspace*{\fill}, which
% can only ever reach the last page; set as the page foot here, so a CV of more
% than one page keeps it.
%
% The bar spans the paper rather than the text block, which takes some care.
% fancyhdr centres the foot with \hfil, and \hfil stretches but does not shrink,
% so simply handing it an overwide box does not centre that box -- the glue sits
% at its natural zero and the box hangs off the right edge alone. The outer
% \makebox is therefore exactly \textwidth, which is the width fancyhdr expects
% and centres properly; the bleed is then applied inside it, where it is
% measured from the left text edge and so reaches the paper edge exactly.
%
% Vertically the bar is lowered onto the bottom edge of the paper, so no white
% shows beneath it. The bar is boxed first because the shift depends on its own
% depth; the \raisebox is given that box's original height and depth so the
% shift moves the bar without also moving the foot it is set in.
\newsavebox{\cvfootbox}
\newcommand{\cvfooterbar}[1]{%
  \sbox\cvfootbox{%
    \colorbox{bgcol}{\makebox[\dimexpr\paperwidth-2\fboxsep\relax][c]{%
      \mystrut\footnotesize\textcolor{white}{#1}}}}%
  \makebox[\textwidth][l]{%
    \hspace*{-\cvbleed}%
    \raisebox{-\dimexpr\cvfootbleed-\dp\cvfootbox\relax}%
      [\ht\cvfootbox][\dp\cvfootbox]{\usebox\cvfootbox}}}

% `footer` fills the bar when it is given. Failing that it carries the web and
% code addresses, which is what upstream puts there -- so the bar is present on
% a document that never asked for one, as it is in the design.
$if(footer)$
\cfoot{\cvfooterbar{$footer$}}
$else$
$if(www)$
\cfoot{\cvfooterbar{\href{https://$www$}{$www$}%
  $if(github)$\cvcontactsep\href{https://github.com/$github$}{github.com/$github$}$endif$}}
$else$
$if(github)$
\cfoot{\cvfooterbar{\href{https://github.com/$github$}{github.com/$github$}}}
$endif$
$endif$
$endif$

\begin{document}

\begin{multicols}{2}

%-------------------------------------------------------------------------------
%	TITLE
%-------------------------------------------------------------------------------
% Set as the first material *inside* the columns, so that it opens the left one
% and the body flows on beneath it. Upstream places it in a hand-built left
% minipage, which a single Markdown stream has no way to address; being first in
% the multicol stream puts it in the same place without needing to.
%
% The bar and the photograph under it are drawn \cvbleed wider than the column
% and pulled that far left, so they run off the page edge as upstream's do.
% \fboxsep is subtracted rather than zeroed, so the padding keeping the name off
% the edge of its box survives and the bar still ends flush with the column.
%
% The accent rule falls between the given and family names, which is what
% \cvtitlerule is for and where upstream puts it -- not beside `docname`, which
% this design has no room for and upstream does not set.
\noindent\hspace*{-\cvbleed}%
\colorbox{bgcol}{\makebox[\dimexpr\linewidth+\cvbleed-2\fboxsep\relax][c]{%
  \HUGE\textcolor{white}{\textsc{$name$}}%
  $if(surname)$\cvtitlerule\textcolor{white}{\textsc{$surname$}}$endif$}}%
$if(profilepic)$
% The photograph sits flush under the bar, with nothing between the two boxes:
% \nointerlineskip drops the interline glue, and \parskip is zeroed for the
% paragraph the picture opens, since the picture is a continuation of the bar
% rather than a paragraph after it.
\par\nointerlineskip
\setlength{\cvbleedwidth}{\dimexpr\linewidth+\cvbleed\relax}%
{\parskip=0pt\relax
 \noindent\hspace*{-\cvbleed}\includegraphics[width=\cvbleedwidth]{$profilepic$}\par}%
$endif$
\par\vspace{8pt}

\normalsize

$latexcv-body.tex()$

\end{multicols}

\end{document}
